\documentclass[12pt,a4paper]{article}
\usepackage[UTF8]{ctex}
\usepackage{amsmath,amssymb,amsthm}
\usepackage{geometry}

\geometry{margin=2.5cm}

\title{为什么阻抗控制的机器人应该只与导纳型环境耦合？}
\author{第11章补充说明}
\date{}

\begin{document}

\maketitle

\section{问题提出}

在机器人控制中，有这样的表述：

"Theoretically, an impedance-controlled robot should only be coupled to an admittance-type environment."

\textbf{问题}：
\begin{enumerate}
\item 什么是阻抗型环境和导纳型环境？
\item 为什么阻抗控制的机器人应该只与导纳型环境耦合？
\item 如果与阻抗型环境耦合会发生什么？
\end{enumerate}

\section{基本概念回顾}

\subsection{阻抗控制}

\textbf{定义}：阻抗控制感测运动，产生力。

\textbf{传递函数}：
\begin{equation}
Z(s) = \frac{F(s)}{X(s)} \quad \text{（阻抗）}
\end{equation}

\textbf{物理意义}：
\begin{itemize}
\item 输入：运动$X$（位置/速度）
\item 输出：力$F$
\item 机器人感测自己的运动，产生相应的力
\end{itemize}

\textbf{实现方式}：
\begin{itemize}
\item 使用编码器、转速计感测运动
\item 根据运动计算并产生力
\item 不需要力传感器（通常）
\end{itemize}

\subsection{导纳控制}

\textbf{定义}：导纳控制感测力，产生运动。

\textbf{传递函数}：
\begin{equation}
Y(s) = \frac{X(s)}{F(s)} = \frac{1}{Z(s)} \quad \text{（导纳）}
\end{equation}

\textbf{物理意义}：
\begin{itemize}
\item 输入：力$F$
\item 输出：运动$X$（位置/速度）
\item 机器人感测外力，产生相应的运动
\end{itemize}

\textbf{实现方式}：
\begin{itemize}
\item 使用力-力矩传感器感测外力
\item 根据外力计算并产生运动
\item 需要力传感器
\end{itemize}

\section{环境类型}

\subsection{阻抗型环境}

\textbf{定义}：环境呈现阻抗特性，感测运动，产生力。

\textbf{传递函数}：
\begin{equation}
Z_{\text{env}}(s) = \frac{F_{\text{env}}(s)}{X(s)} \quad \text{（环境阻抗）}
\end{equation}

\textbf{物理意义}：
\begin{itemize}
\item 环境感测运动$X$（如位置变化）
\item 环境产生力$F_{\text{env}}$（如弹簧力、阻尼力）
\item 例如：弹簧、阻尼器、刚性墙
\end{itemize}

\textbf{例子}：
\begin{itemize}
\item \textbf{弹簧}：$F = kx$（感测位置，产生力）
\item \textbf{阻尼器}：$F = b\dot{x}$（感测速度，产生力）
\item \textbf{刚性墙}：感测穿透，产生反作用力
\end{itemize}

\subsection{导纳型环境}

\textbf{定义}：环境呈现导纳特性，感测力，产生运动。

\textbf{传递函数}：
\begin{equation}
Y_{\text{env}}(s) = \frac{X(s)}{F(s)} \quad \text{（环境导纳）}
\end{equation}

\textbf{物理意义}：
\begin{itemize}
\item 环境感测力$F$（如外力）
\item 环境产生运动$X$（如位移）
\item 例如：自由空间、可移动物体、另一个机器人
\end{itemize}

\textbf{例子}：
\begin{itemize}
\item \textbf{自由空间}：感测力，产生运动（几乎无阻力）
\item \textbf{可移动物体}：感测推力，产生位移
\item \textbf{另一个导纳控制的机器人}：感测力，产生运动
\end{itemize}

\section{为什么阻抗控制应该与导纳型环境耦合？}

\subsection{基本原理}

\textbf{关键原则}：两个系统要能够耦合，它们的输入输出必须匹配。

\textbf{匹配规则}：
\begin{itemize}
\item 一个系统的输出应该是另一个系统的输入
\item 一个系统的输入应该是另一个系统的输出
\end{itemize}

\subsection{阻抗控制 + 导纳型环境}

\textbf{阻抗控制的机器人}：
\begin{itemize}
\item 输入：运动$X$
\item 输出：力$F_{\text{robot}}$
\end{itemize}

\textbf{导纳型环境}：
\begin{itemize}
\item 输入：力$F$
\item 输出：运动$X_{\text{env}}$
\end{itemize}

\textbf{耦合关系}：
\begin{itemize}
\item 机器人的输出（力$F_{\text{robot}}$）是环境的输入
\item 环境的输出（运动$X_{\text{env}}$）是机器人的输入
\item \textcolor{green}{完美匹配！}
\end{itemize}

\textbf{框图}：
\begin{center}
机器人（阻抗）$\to$ 力$F$ $\to$ 环境（导纳）$\to$ 运动$X$ $\to$ 机器人（阻抗）
\end{center}

\textbf{物理意义}：
\begin{itemize}
\item 机器人感测运动，产生力
\item 环境感测力，产生运动
\item 两者形成闭环，可以稳定工作
\end{itemize}

\subsection{阻抗控制 + 阻抗型环境（不匹配）}

\textbf{阻抗控制的机器人}：
\begin{itemize}
\item 输入：运动$X$
\item 输出：力$F_{\text{robot}}$
\end{itemize}

\textbf{阻抗型环境}：
\begin{itemize}
\item 输入：运动$X$
\item 输出：力$F_{\text{env}}$
\end{itemize}

\textbf{问题}：
\begin{itemize}
\item 两者都需要运动作为输入
\item 两者都产生力作为输出
\item \textcolor{red}{输入输出不匹配！}
\end{itemize}

\textbf{冲突}：
\begin{itemize}
\item 机器人感测运动，产生力
\item 环境也感测运动，产生力
\item 两者都试图控制力，可能产生冲突
\item 可能导致不稳定或不可预测的行为
\end{itemize}

\section{详细分析}

\subsection{阻抗控制 + 导纳型环境（匹配）}

\textbf{系统方程}：

\textbf{机器人（阻抗控制）}：
\begin{equation}
F_{\text{robot}} = Z_{\text{robot}}(s) X(s)
\end{equation}

\textbf{环境（导纳型）}：
\begin{equation}
X(s) = Y_{\text{env}}(s) F(s)
\end{equation}

\textbf{耦合关系}：
\begin{equation}
F(s) = F_{\text{robot}}(s) \quad \text{（机器人的力作用于环境）}
\end{equation}

\textbf{闭环系统}：
\begin{align}
F_{\text{robot}} &= Z_{\text{robot}} X \\
X &= Y_{\text{env}} F_{\text{robot}} \\
X &= Y_{\text{env}} Z_{\text{robot}} X
\end{align}

\textbf{稳定性}：
\begin{itemize}
\item 闭环传递函数：$G_{\text{closed}}(s) = Y_{\text{env}}(s) Z_{\text{robot}}(s)$
\item 如果$Y_{\text{env}}$和$Z_{\text{robot}}$都是稳定的，闭环系统稳定
\item 两者是互补的，可以形成稳定的闭环
\end{itemize}

\subsection{阻抗控制 + 阻抗型环境（不匹配）}

\textbf{系统方程}：

\textbf{机器人（阻抗控制）}：
\begin{equation}
F_{\text{robot}} = Z_{\text{robot}}(s) X(s)
\end{equation}

\textbf{环境（阻抗型）}：
\begin{equation}
F_{\text{env}} = Z_{\text{env}}(s) X(s)
\end{equation}

\textbf{问题}：
\begin{itemize}
\item 两者都需要$X$作为输入
\item 两者都产生力作为输出
\item 没有形成闭环
\end{itemize}

\textbf{力冲突}：
\begin{itemize}
\item 机器人根据运动产生力$F_{\text{robot}}$
\item 环境根据运动产生力$F_{\text{env}}$
\item 总力：$F_{\text{total}} = F_{\text{robot}} + F_{\text{env}}$
\item 两者可能冲突，导致不稳定
\end{itemize}

\textbf{不稳定示例}：

假设：
\begin{itemize}
\item 机器人：$F_{\text{robot}} = K X$（高刚度）
\item 环境：$F_{\text{env}} = k X$（高刚度）
\end{itemize}

如果两者都试图控制位置，可能产生振荡或不稳定。

\section{对偶关系}

\subsection{阻抗 vs 导纳}

\textbf{对偶性}：
\begin{itemize}
\item 阻抗$Z$和导纳$Y$是互逆的：$Y = 1/Z$
\item 阻抗控制感测运动，产生力
\item 导纳控制感测力，产生运动
\item 两者是对偶的
\end{itemize}

\textbf{匹配规则}：
\begin{itemize}
\item 阻抗控制的机器人 $\leftrightarrow$ 导纳型环境
\item 导纳控制的机器人 $\leftrightarrow$ 阻抗型环境
\end{itemize}

\subsection{四种组合}

\begin{center}
\begin{tabular}{|l|l|l|}
\hline
\textbf{机器人类型} & \textbf{环境类型} & \textbf{是否匹配} \\
\hline
阻抗控制 & 导纳型 & \textcolor{green}{是} \\
\hline
阻抗控制 & 阻抗型 & \textcolor{red}{否} \\
\hline
导纳控制 & 阻抗型 & \textcolor{green}{是} \\
\hline
导纳控制 & 导纳型 & \textcolor{red}{否} \\
\hline
\end{tabular}
\end{center}

\section{实际例子}

\subsection{例子1：阻抗控制机器人 + 自由空间（导纳型）}

\textbf{场景}：机器人在自由空间中运动

\textbf{机器人}：
\begin{itemize}
\item 阻抗控制：感测运动，产生力
\item 在自由空间中，运动不受限制
\end{itemize}

\textbf{环境}：
\begin{itemize}
\item 自由空间：导纳型（感测力，产生运动）
\item 几乎无阻力，力产生运动
\end{itemize}

\textbf{结果}：\textcolor{green}{匹配，可以正常工作}

\subsection{例子2：阻抗控制机器人 + 刚性墙（阻抗型）}

\textbf{场景}：机器人接触刚性墙

\textbf{机器人}：
\begin{itemize}
\item 阻抗控制：感测运动，产生力
\item 试图根据位置产生力
\end{itemize}

\textbf{环境}：
\begin{itemize}
\item 刚性墙：阻抗型（感测穿透，产生反作用力）
\item 也根据位置产生力
\end{itemize}

\textbf{问题}：
\begin{itemize}
\item 两者都试图控制力
\item 可能产生冲突
\item 可能导致振荡或不稳定
\end{itemize}

\textbf{结果}：\textcolor{red}{不匹配，可能不稳定}

\subsection{例子3：导纳控制机器人 + 刚性墙（阻抗型）}

\textbf{场景}：导纳控制机器人接触刚性墙

\textbf{机器人}：
\begin{itemize}
\item 导纳控制：感测力，产生运动
\item 感测接触力，调整运动
\end{itemize}

\textbf{环境}：
\begin{itemize}
\item 刚性墙：阻抗型（感测穿透，产生反作用力）
\item 根据位置产生力
\end{itemize}

\textbf{耦合关系}：
\begin{itemize}
\item 环境感测位置，产生力
\item 机器人感测力，产生运动
\item \textcolor{green}{完美匹配！}
\end{itemize}

\textbf{结果}：\textcolor{green}{匹配，可以正常工作}

\section{为什么理论上是这样？}

\subsection{理论依据}

\textbf{系统理论}：
\begin{itemize}
\item 两个系统要形成稳定的闭环，输入输出必须匹配
\item 一个系统的输出应该是另一个系统的输入
\item 否则无法形成闭环，或形成不稳定的闭环
\end{itemize}

\textbf{控制理论}：
\begin{itemize}
\item 阻抗控制和导纳控制是对偶的
\item 对偶系统可以形成稳定的闭环
\item 非对偶系统可能产生冲突
\end{itemize}

\subsection{实际考虑}

\textbf{理论 vs 实际}：
\begin{itemize}
\item 理论上：阻抗控制应该只与导纳型环境耦合
\item 实际上：由于柔顺性、延迟等因素，可能可以与其他类型环境工作
\item 但性能可能不是最优的
\end{itemize}

\textbf{为什么说"理论上"？}
\begin{itemize}
\item 在实际系统中，总是存在一些柔顺性
\item 关节和连杆的柔顺性提供了"缓冲"
\item 这使得阻抗控制也可以与某些阻抗型环境工作
\item 但这不是最优配置
\end{itemize}

\section{总结}

\subsection{核心原理}

\begin{enumerate}
\item \textbf{匹配规则}：
\begin{itemize}
\item 阻抗控制的机器人（感测运动，产生力）应该与导纳型环境（感测力，产生运动）耦合
\item 导纳控制的机器人（感测力，产生运动）应该与阻抗型环境（感测运动，产生力）耦合
\end{itemize}

\item \textbf{原因}：
\begin{itemize}
\item 输入输出匹配：一个系统的输出是另一个系统的输入
\item 形成稳定的闭环
\item 避免力冲突
\end{itemize}

\item \textbf{不匹配的后果}：
\begin{itemize}
\item 输入输出不匹配
\item 可能产生力冲突
\item 可能导致不稳定
\end{itemize}
\end{enumerate}

\subsection{关键公式}

\begin{enumerate}
\item \textbf{阻抗控制}：$F = Z(s) X$
\item \textbf{导纳控制}：$X = Y(s) F$
\item \textbf{匹配关系}：$Y = 1/Z$
\item \textbf{闭环稳定性}：需要输入输出匹配
\end{enumerate}

\subsection{实际应用}

\begin{itemize}
\item \textbf{设计机器人时}：考虑环境类型，选择合适的控制方法
\item \textbf{选择控制方法时}：考虑环境的阻抗/导纳特性
\item \textbf{实际实现时}：由于柔顺性，可能可以放宽理论限制，但要注意稳定性
\end{itemize}

\end{document}

